%% [2026-07-27] NEW SECTION -- entirely proposed, typeset in red via \rev{}.
%% Replaces the "Literature." paragraph currently sitting at the end of the introduction,
%% which should be deleted once this section is adopted (see the note reported in chat).
%%
%% Citation protocol for this section:
%%   UPDATED  gj2023global -- was "ECB Working Paper 2881" (2023); now published as
%%            Journal of International Economics 158, 104169 (2025).
%%   ADDED    boeck2025globalization (Boeck & Mori, JIE 157, 2025)
%%            boeck2021fg            (Boeck, Feldkircher & Siklos, OBES 83(5), 2021)
%%            boeck2024carry         (Boeck, Steshkova & Zoerner, OeNB WP 258, 2024)
%%   DROPPED  ca2020monetary  -- ECB WP superseded by the published ca2021making.
%%            kim2001international, canova2005transmission -- pre-crisis, superseded on
%%            the same point by the post-2020 references retained below.
%%            dees2007exploring, feldkircher2016international, koop2016model,
%%            burriel2018uncovering -- the GVAR cluster is now represented by
%%            georgiadis2017bi and dees2021global rather than by a six-item list.
%%   RETAINED bluwstein2018beggar, crespo2019spillovers, breitenlechner2021spillbacks,
%%            ca2021making, dees2021global, georgiadis2017bi, miranda2020us,
%%            gopinath2021dominant, tenreyro2016pushing, altavilla2019measuring,
%%            jarocinski2020deconstructing, DIEBOLD2006337.

\section{Related Literature}\label{sec:literature}

\rev{Our question sits at the intersection of three strands of work. Each supplies one ingredient of the exercise, but none combines them in the way required to ask whether the effect of a domestic policy surprise depends on the stance of a foreign central bank.}

\rev{The first strand documents the \textit{international transmission of US monetary policy}. \cite{miranda2020us} show that a global financial cycle in asset prices, leverage and capital flows is driven in large part by the Fed, operating through risk premia rather than through interest parity alone, while \cite{gopinath2021dominant} trace a complementary channel running through the dollar's role as the dominant invoicing and funding currency. Recent contributions have sharpened both the shock and the mapping. \citet{gj2023global} decompose US policy into conventional, forward-guidance and asset-purchase dimensions and find that the latter two generate the larger spillovers, transmitted predominantly through financial channels. \citet{breitenlechner2021spillbacks} show that these spillovers return to their source as quantitatively relevant spillbacks, and \citet{ca2021making} document that transmission in a globalized world is asymmetric across countries and instruments. \citet{boeck2025globalization} ask whether the international transmission of US policy has itself changed as globalization proceeded, and find trade channels rivalling and eventually surpassing financial ones. Spillovers also run in the opposite direction, if more weakly: \citet{boeck2021fg} document international effects of euro area forward guidance, and \citet{dees2021global} place the global financial cycle in an interconnected multi-country setting. What unites this literature is that the \textit{foreign} policy shock is the impulse and the domestic economy the recipient. We invert these roles. The Fed's stance is not our shock but our conditioning variable, and the impulse is a domestic one.}

\rev{The second strand studies \textit{state dependence in monetary transmission}: whether a given policy shock propagates differently depending on the environment into which it is delivered. The bulk of this work conditions on \textit{domestic} conditions, most prominently the business cycle, and reaches mixed conclusions \citep[see, e.g.,][]{tenreyro2016pushing}. Closer to our design is work conditioning on financial-market states: \citet{boeck2024carry} use a threshold VAR to show that the transmission of policy to currency markets depends on the intensity of carry trade activity, and \citet{bluwstein2018beggar} and \citet{crespo2019spillovers} document that the size of cross-border effects varies with the financial environment. This literature establishes that transmission is state dependent and provides the econometric apparatus for measuring it, but the state is almost always domestic. Our conditioning variable is the policy stance of a foreign central bank, which is conceptually distinct from domestic slack and, over our sample, moves largely independently of it.}

\rev{The third strand supplies the measurement. High-frequency identification recovers policy surprises from market movements in narrow windows around announcements, and \citet{altavilla2019measuring} provide the euro area event-study database and factor decomposition we use. Because such surprises conflate policy news with information about the central bank's outlook, \citet{jarocinski2020deconstructing} propose a sign-restriction scheme that separates the two; we adopt their correction and use their shocks as an alternative identification in our robustness exercise. To move from individual maturities to the curve, we rely on the Nelson-Siegel tradition in the form of \citet{DIEBOLD2006337}, which lets us summarize fifteen maturities of nominal yields and inflation-linked swaps by a small set of factors and recover responses across the entire term structure.}

\rev{Our contribution is to combine the conditioning variable of the first strand with the question of the second, using the measurement apparatus of the third. Relative to the spillover literature, we do not ask how large the effects of a Fed shock are abroad, but whether the Fed's prevailing stance changes the transmission of someone else's policy. Relative to the state-dependence literature, we replace a domestic state with a foreign one, and we do so along a continuum rather than by splitting the sample into two regimes: the smooth transition function delivers a distinct set of dynamic coefficients for every value of the Fed's stance. Relative to both, we report the response of the full term structure of nominal yields, inflation compensation and real rates rather than of a single benchmark maturity. This last choice turns out to matter, because the state dependence we document is invisible at the short end and concentrated at the long end, and because the decomposition into real rates and inflation compensation locates where the composition of the adjustment changes and where it does not. To our knowledge, the interaction between the ECB's transmission and the Fed's stance has not previously been documented along the term structure. \citet{georgiadis2017bi} shows that estimates of cross-border effects depend materially on whether the model is bilateral or multilateral, which is a caution we take seriously in interpreting the magnitude of our own estimates.}
